\subsection{Métricas de evaluación}

El desempeño se evalúa reuniendo las predicciones de un paso obtenidas en todas las ventanas y comparándolas con sus generaciones observadas. Las métricas se calculan sobre estas predicciones posteriores al ajuste de cada ventana. Sea $m$ el número de predicciones evaluadas, $g_i$ el valor observado, $\hat g_i$ la predicción correspondiente y $\bar g$ la media de los valores observados evaluados.

El error absoluto medio (MAE) cuantifica la magnitud promedio de las diferencias:
\begin{equation}
\mathrm{MAE}=\frac{1}{m}\sum_{i=1}^{m}|g_i-\hat g_i|.
\end{equation}
Se expresa en MWh y un valor menor indica mayor precisión.

La raíz del error cuadrático medio (RMSE) se calcula mediante
\begin{equation}
\mathrm{RMSE}=\sqrt{\frac{1}{m}\sum_{i=1}^{m}(g_i-\hat g_i)^2}.
\end{equation}
También se expresa en MWh y otorga mayor peso a las desviaciones grandes.

El coeficiente de determinación compara el error del modelo con la variabilidad de las observaciones:
\begin{equation}
R^2=1-\frac{\sum_{i=1}^{m}(g_i-\hat g_i)^2}
{\sum_{i=1}^{m}(g_i-\bar g)^2}.
\end{equation}
Esta expresión requiere que los valores observados no sean constantes. Un valor de uno corresponde a predicciones exactas; cero indica un error cuadrático igual al de predecir la media del conjunto evaluado, y un valor negativo indica un error mayor.

El programa muestra MAE, RMSE y $R^2$ con dos decimales, junto con el área, el intervalo entre registros, las fechas evaluadas y el número de predicciones. Además, guarda una gráfica de generación observada y estimada frente a fecha y hora. Los resultados numéricos de la ejecución sobre los datos del estudio se incorporarán al capítulo de resultados.

\subsubsection{Error porcentual absoluto medio}

Como medida complementaria se considera el error porcentual absoluto medio (MAPE), cuya expresión aparece en el cuaderno \texttt{ventanas.ipynb}:
\begin{equation}
\mathrm{MAPE}=\frac{100}{m}\sum_{i=1}^{m}
\left|\frac{g_i-\hat g_i}{g_i}\right|.
\end{equation}
Esta métrica expresa el error absoluto relativo promedio como porcentaje y requiere valores observados distintos de cero. Los valores pequeños de generación pueden aumentar su sensibilidad. El programa \texttt{root/OLS.py} calcula MAE, RMSE y $R^2$; MAPE se incorpora aquí como evaluación complementaria del cuaderno y deberá calcularse sobre las mismas predicciones cuando se complete una ejecución válida.
